% ============================================================================
% Fichier  : partie4/P04_C20_nuage_mobiles.tex
% Titre    : Sécurité du Nuage et des Mobiles
% Auteur   : MINKA MI NGUIDJOÏ Thierry Emmanuel
% Version  : 1.0 -- novembre 2026
% ============================================================================
% Sources : NIST SP 800-145 (définition du nuage, 2011) ; CSA Cloud Controls
% Matrix v4 ; ISO/IEC 27017:2015 ; NIST SP 800-124 Rev. 2 (mobiles, 2023) ;
% OWASP MASVS.
% ============================================================================

\chapter{Sécurité du Nuage et des Mobiles}
\label{chap:nuage}
\chaptermark{Nuage et mobiles}

\epigraph{%
    « Le nuage, c'est l'ordinateur de quelqu'un d'autre~;\\
    la configuration, c'est toujours la vôtre. »%
}{Principe du cours}

\niveauChapitre{Intermédiaire}{Chapitres~\ref{chap:architecture},
\ref{chap:iam} et~\ref{chap:applications}.}

\begin{boxobjectifs}
\begin{itemize}
    \item Définir les modèles de service et de déploiement du nuage.
    \item Répartir les responsabilités entre fournisseur et client.
    \item Prévenir les erreurs de configuration les plus fréquentes.
    \item Protéger les identités, les journaux et les clés dans le nuage.
    \item Sécuriser une flotte de téléphones et les applications mobiles.
\end{itemize}
\end{boxobjectifs}

\bigskip

% ============================================================================
\section*{Question directrice : la messagerie de LiZbiZ était-elle à l'abri parce qu'elle était dans le nuage~?}
% ============================================================================

La messagerie de LiZbiZ, hébergée par un grand fournisseur, n'a pas été
chiffrée en mars. La direction en conclut que « le nuage est sûr ». Une
vérification montre pourtant que les comptes d'administration de la
messagerie n'ont pas de second facteur et que les journaux d'accès ne
sont pas conservés. Le fournisseur protège son infrastructure~; la
configuration et les comptes restent l'affaire du client.

\medskip
\textbf{Plan du chapitre.}
\begin{enumerate}
    \item Les modèles du nuage (section~\ref{sec:c20-modeles}).
    \item La responsabilité partagée (section~\ref{sec:c20-partage}).
    \item Les erreurs de configuration (section~\ref{sec:c20-config}).
    \item Identités, journaux et clés (section~\ref{sec:c20-identites}).
    \item Les mobiles (section~\ref{sec:c20-mobiles}).
\end{enumerate}

\separateur

% ============================================================================
\section{Les Modèles du Nuage}
\label{sec:c20-modeles}
% ============================================================================

Le NIST définit l'informatique en nuage par cinq caractéristiques
(libre-service à la demande, accès par le réseau, mutualisation des
ressources, élasticité, service mesuré), trois modèles de service et
quatre modèles de déploiement~\cite{nist2011sp800145}.

\begin{tableaucours}{L{2.4cm}L{5.0cm}L{4.8cm}}{%
    \tableauHeader{Modèle} &
    \tableauHeader{Le fournisseur livre} &
    \tableauHeader{Exemple}
}
IaaS & Machines virtuelles, stockage, réseau & Serveurs loués \\
PaaS & Plateforme d'exécution, bases gérées & Base de données gérée \\
SaaS & Application complète & Messagerie, paie en ligne \\
\end{tableaucours}
\vspace{-0.8em}
\begin{center}{\small\itshape Tableau~20.1. Les modèles de service}\end{center}

Les modèles de déploiement sont le nuage public, privé, communautaire et
hybride.

% ============================================================================
\section{La Responsabilité Partagée}
\label{sec:c20-partage}
% ============================================================================

\begin{tableaucours}{L{4.0cm}ccc}{%
    \tableauHeader{Couche} &
    \tableauHeader{IaaS} &
    \tableauHeader{PaaS} &
    \tableauHeader{SaaS}
}
Données et comptes & Client & Client & Client \\
Configuration de l'application & Client & Client & Client \\
Application & Client & Client & Fournisseur \\
Système d'exploitation & Client & Fournisseur & Fournisseur \\
Infrastructure physique & Fournisseur & Fournisseur & Fournisseur \\
\end{tableaucours}
\vspace{-0.8em}
\begin{center}{\small\itshape Tableau~20.2. Qui protège quoi selon le modèle}\end{center}

Dans tous les modèles, le client reste responsable de ses données, de ses
comptes et de sa configuration. La CSA décrit les contrôles à répartir
entre les deux parties dans sa Cloud Controls Matrix~\cite{csa2021ccm}~;
l'ISO/IEC~27017 complète l'ISO/IEC~27002 pour le nuage~\cite{iso27017_2015}.

% ============================================================================
\section{Les Erreurs de Configuration}
\label{sec:c20-config}
% ============================================================================

\begin{tableaucours}{L{4.4cm}L{7.8cm}}{%
    \tableauHeader{Erreur fréquente} &
    \tableauHeader{Mesure}
}
Stockage d'objets ouvert au public & Interdire l'accès public par défaut au
niveau du compte \\
Comptes d'administration sans second facteur & Second facteur obligatoire,
comptes d'urgence protégés \\
Clés d'accès dans le code ou anciennes & Coffre de secrets, rotation \\
Groupes de sécurité ouverts à tout Internet & Filtrer par adresse et par
port \\
Journaux désactivés ou trop courts & Activer et exporter les journaux \\
Ressources oubliées & Inventaire, étiquettes, revue des coûts \\
\end{tableaucours}
\vspace{-0.8em}
\begin{center}{\small\itshape Tableau~20.3. Erreurs de configuration et mesures}\end{center}

\begin{boxalert}
\textbf{Violation typique.} Rendre un espace de stockage public « pour
partager un fichier avec un client », puis oublier.

\smallskip
\textbf{Règle~:} on partage par un lien limité dans le temps ou par un
accès nominatif~; l'accès public est bloqué par défaut au niveau du
compte, et toute exception est revue.
\end{boxalert}

% ============================================================================
\section{Identités, Journaux et Clés}
\label{sec:c20-identites}
% ============================================================================

Dans le nuage, le périmètre est l'identité~: qui tient un compte
d'administration de la console tient l'infrastructure. Les journaux
d'activité du fournisseur (connexions, changements de configuration,
appels d'interface) doivent être activés, exportés et conservés hors du
compte qu'ils surveillent. Pour les données les plus sensibles, le client
peut garder la maîtrise des clés de chiffrement (clés gérées par le
client), au prix d'une gestion plus lourde.

\begin{boxscenario}{Opération GAROUA. La messagerie vérifiée}
\textbf{Constatations.} Trois comptes d'administration de la messagerie,
sans second facteur~; journaux d'accès conservés 7~jours seulement, sans
export~; règle de transfert automatique, créée en février, qui envoie à
une adresse externe les courriels contenant le mot « facture ».

\medskip
\textbf{Analyse} (confiance moyenne). La règle de transfert est un signe
classique de compromission d'une boîte~; sa création en février coïncide
avec l'hameçonnage. Faute de journaux au-delà de 7~jours, on ne peut pas
établir qui l'a créée~: la question reste ouverte (\NEK{}).
\end{boxscenario}

% ============================================================================
\section{Les Mobiles}
\label{sec:c20-mobiles}
% ============================================================================

Les téléphones portent la messagerie, les codes d'authentification et
souvent le Mobile Money~: ce sont des postes de travail à part entière. Le
NIST décrit leur gestion en entreprise~\cite{nist800124}.

\begin{tableaucours}{L{4.0cm}L{8.2cm}}{%
    \tableauHeader{Mesure} &
    \tableauHeader{Effet}
}
Gestion de flotte (MDM) & Inventaire, configuration imposée, effacement à
distance \\
Séparation professionnel et personnel & Conteneur professionnel sur les
appareils personnels (BYOD) \\
Verrouillage et chiffrement & Protègent les données en cas de perte \\
Applications depuis les magasins officiels & Limitent les applications
piégées \\
Mises à jour du système & Corrigent les failles exploitées \\
\end{tableaucours}
\vspace{-0.8em}
\begin{center}{\small\itshape Tableau~20.4. Sécuriser une flotte mobile}\end{center}

Pour les développeurs d'applications mobiles, l'OWASP publie un référentiel
d'exigences, le MASVS~: stockage des données, cryptographie,
authentification, communications, protection du code~\cite{owaspmasvs}.

\begin{boxmetier}{RSSI}
\textbf{Ne pas déléguer sa vigilance.} Le RSSI exige du fournisseur ses
rapports d'assurance, mais vérifie lui-même ce qui relève du client~:
comptes, configuration, journaux, partages.
\end{boxmetier}

\begin{boxmetier}{Analyste SOC}
\textbf{Surveiller la console.} L'analyste alerte sur les connexions
d'administration inhabituelles, la création de règles de transfert, la
désactivation des journaux et l'ouverture d'un stockage au public.
\end{boxmetier}

\separateur

\begin{boxresume}
\textbf{Ce qu'il faut retenir de ce chapitre~:}
\begin{itemize}
    \item \textbf{Le nuage} se décline en IaaS, PaaS et SaaS, et en
          déploiements public, privé, communautaire, hybride.
    \item \textbf{La responsabilité est partagée}~: données, comptes et
          configuration restent toujours au client.
    \item \textbf{Les erreurs de configuration} (stockage public, comptes
          sans second facteur, journaux absents) sont la première cause
          d'incident.
    \item \textbf{L'identité est le périmètre}~; les journaux s'exportent
          hors du compte surveillé.
    \item \textbf{Les mobiles} sont des postes de travail~: gestion de
          flotte, séparation, chiffrement, mises à jour~; MASVS pour les
          applications.
\end{itemize}
\end{boxresume}

\begin{boxinfo}
\textbf{Points de contrôle de l'élève.} Je sais~: définir IaaS, PaaS et
SaaS~; remplir le tableau de responsabilité partagée~; citer cinq erreurs
de configuration~; dire pourquoi exporter les journaux~; proposer les
mesures d'une flotte mobile.
\end{boxinfo}

\bigskip

\begin{boxexo}[title={\ding{45}\quad Exercice~20.1. Quel modèle~? \niveaubase}]
Classez~: une messagerie en ligne~; une machine virtuelle louée~; une base
de données gérée par le fournisseur.
\end{boxexo}

\begin{boxexo}[title={\ding{45}\quad Exercice~20.2. Qui est responsable~? \niveaubase}]
Pour une messagerie en SaaS, qui est responsable~: du correctif du serveur~;
du second facteur des administrateurs~; des règles de transfert~; de la
sauvegarde des courriels~?
\end{boxexo}

\begin{boxexo}[title={\ding{45}\quad Exercice~20.3. Le stockage public \niveaubase}]
Pourquoi un stockage d'objets public est-il si dangereux~? Quelle mesure
au niveau du compte l'empêche~?
\end{boxexo}

\begin{boxexo}[title={\ding{45}\quad Exercice~20.4. La règle de transfert \niveaumoyen}]
Rédigez séparément la constatation et l'analyse sur la règle de transfert
de la messagerie de LiZbiZ, et les actions à mener.
\end{boxexo}

\begin{boxexo}[title={\ding{45}\quad Exercice~20.5. Les journaux \niveaumoyen}]
Proposez une configuration des journaux de la messagerie~: événements,
export, durée de conservation, alertes.
\end{boxexo}

\begin{boxexo}[title={\ding{45}\quad Exercice~20.6. BYOD \niveaumoyen}]
Les commerciaux de LiZbiZ utilisent leurs téléphones personnels.
Proposez une politique qui protège LiZbiZ sans surveiller la vie privée
des salariés.
\end{boxexo}

\begin{boxexo}[title={\ding{45}\quad Exercice~20.7. Clés \niveaumoyen}]
Quels avantages et quels risques présente la gestion des clés de
chiffrement par le client~?
\end{boxexo}

\begin{boxexo}[title={\ding{45}\quad Exercice~20.8. Un téléphone perdu \niveaumoyen}]
Le téléphone du directeur financier, qui porte son second facteur et
l'accès au Mobile Money de l'entreprise, est perdu. Rédigez la procédure.
\end{boxexo}

\begin{boxexo}[title={\ding{45}\quad Exercice~20.9. Migrer \niveauexpert}]
LiZbiZ veut déplacer son application intégrée vers un IaaS. Rédigez les
dix exigences de sécurité à intégrer au projet.
\end{boxexo}

\begin{boxexo}[title={\ding{45}\quad Exercice~20.10. Revue de configuration \niveauexpert}]
Construisez une liste de contrôle de vingt points pour revoir la
configuration d'un compte SaaS de messagerie, avec la preuve attendue
pour chacun.
\end{boxexo}

\bigskip

\begin{boxinfo}
\textbf{Pour aller plus loin.}
\begin{itemize}
    \item NIST SP~800-145 et SP~800-124~\cite{nist2011sp800145,nist800124}.
    \item CSA Cloud Controls Matrix~\cite{csa2021ccm}~; ISO/IEC~27017~\cite{iso27017_2015}.
    \item OWASP MASVS~\cite{owaspmasvs}.
\end{itemize}
\end{boxinfo}

\begin{boxinfo}
\textbf{Transition.} Les mesures sont en place~; encore faut-il vérifier
qu'elles fonctionnent. Le Chapitre~\ref{chap:evaluer} traite de
l'évaluation de la sécurité.
\end{boxinfo}
